%%%%%%%%%%%%%%%%%%%%%%%%%%%%%%%%%
%
%       EXERCÍCIO
%
%%%%%%%%%%%%%%%%%%%%%%%%%%%%%%%%%

\definevalorquestao

\begin{exercicioBanco}[\valorquestao]
Considere as seguintes afirmações sobre exponenciação e suas propriedades:

\begin{enumerate}[label=\Roman*.]
    \item O valor da expressão \((-4)^{0{,}5}\) é igual a \(-2\).
    \item A função \(g(x) = \left(\dfrac{1}{\pi}\right)^x\) é decrescente para todo \(x \in \mathbb{R}\).
    \item Para todo \(x \in \mathbb{R}\), a igualdade \(2^{x+3} - 2^x = 7 \cdot 2^x\) é verdadeira.
    \item Se \(f(x) = a^x\) é uma função exponencial, então \(f(x+y) = f(x) \cdot f(y)\) para quaisquer \(x, y \in \mathbb{R}\).
\end{enumerate}

Assinale a alternativa correta:

% Define as alternativas
\newcommand{\alternativas}{%
    \begin{center}
        \begin{tabularx}{\textwidth}{XXX}
            (a) Todas são falsas. &
            (b) Todas são verdadeiras. &
            (c) Apenas I e III são falsas. \\[5pt]
            (d) Apenas a afirmação I é falsa. &
            (e) Apenas II e IV são verdadeiras.
        \end{tabularx}
    \end{center}
}

% Define a resposta correta
\newcommand{\resposta}{D}

% Lógica condicional para exibição
\ifdefstring{\atividade}{lista}{%
    \alternativas
    \vspace{0.5em}
    
    \noindent\textbf{Resposta:} letra \textbf{\resposta}.
}{%
    \ifdefstring{\modo}{objetivo}{%
        \alternativas
    }{%
        % modo = discursiva → não mostra alternativas
    }
}
\end{exercicioBanco}

